\subsection{域上的赋值}

命题 2. 设 K 是一个（不必交换的）域，v 是 K 的以 Γ 为值群的赋值。则：

\begin{enumerate}
\def\labelenumi{(\roman{enumi})}
\item
  对 \(x \neq 0\)，\(v(x) \neq +\infty\)。
\item
  集 \(\mathbf{A}\) 由 \(x\in \mathbf{K}\) 中使 \(v(x)\geqslant 0\) 的元素组成，它是 \(\mathbf{K}\) 的子环。
\item
  对所有 \(\alpha \geqslant 0\)（在 \(\Gamma\) 中），集 \(\mathrm{V}_{\alpha}\)（相应地，\(\mathrm{V}_{\alpha}^{\prime}\)）由 \(x \in \mathbf{A}\) 中使 \(v(x) > \alpha\)（相应地，\(v(x) \geqslant \alpha\)）的元素组成，它是 \(\mathbf{A}\) 的双边理想，且每个（左或右）理想 \(\neq (0)\)（\(\mathbf{A}\) 的）都包含某个 \(\mathrm{V}_{\alpha}^{\prime}\)。
\item
  集 \(\mathfrak{m}(\mathbf{A})\) 由 \(x \in \mathbf{A}\) 中使 \(v(x) > 0\) 的元素组成，它是 \(\neq \mathbf{A}\) 的最大理想（在 \(\mathbf{A}\) 中）；\(\mathbf{U}(\mathbf{A}) = \mathbf{A} - \mathfrak{m}(\mathbf{A})\) 是 \(\mathbf{A}\) 的可逆元素所成之集，而 \(\kappa(\mathbf{A}) = \mathbf{A} / \mathfrak{m}(\mathbf{A})\) 是一个（不必交换的）域。
\item
  对所有 \(x \in \mathbf{K} - \mathbf{A}\)，\(x^{-1} \in \mathfrak{m}(\mathbf{A})\)。
\end{enumerate}

断言 (i) 由 \(\overline{v}^{-1}(+\infty)\) 是 K 的一个不等于 K 的理想这一事实得出。A 是环以及诸 \(V_{\alpha}\)、\(V_{\alpha}^{\prime}\) 是双边理想这一事实的验证，凭公理（\(VL_{I}\)）、（\(VL_{II}\)）与（\(VL_{III}\)）是平凡的。若 a 是 A 的一个（例如左）理想，且 \(x \neq 0\) 属于 A，则每个 \(y \in A\)，只要 \(v(y) \geqslant v(x)\)，都可写成 y = zx，其中 \(z = yx^{-1}\)，于是 \(v(z) = v(y) - v(x) \geqslant 0\)，从而 \(z \in A\)；换言之，左理想 Ax 包含诸 \(V_{\alpha}^{\prime}\)（对 \(\alpha \geqslant v(x)\)）。集 \(U(A) = A - m(A)\) 是由 \(x \in K\) 中使 \(v(x) = 0\) 的元素组成之集；若 \(x \in U(A)\)，则

\[
v (x ^ {- 1}) = - v (x) = 0,
\]

从而 \(x^{-1} \in \mathrm{U}(\mathrm{A})\)；反之，若 \(y \in \mathrm{A}\) 在 \(\mathrm{A}\) 中可逆，则 \(v(y) \geqslant 0\)，\(v(y^{-1}) \geqslant 0\)，且 \(v(y) + v(y^{-1}) = 0\)，于是 \(v(y) = 0\)，即 \(y \in \mathrm{U}(\mathrm{A})\)；这证明了 (iv)，而 (v) 由定义立即得出。

A（相应地，m(A)、κ(A)）称为 K 上赋值 v 的环（相应地，理想、剩余域）。

显然 U(A) 是同态 \(v: K^{*} \to \Gamma\) 的核，且像 \(v(\mathbf{K}^{*})\)（乘法群 \(K^{*}\) 在 v 之下的）是加法群 \(\Gamma\) 的子群，称为 v 的阶群或值群，它因此同构于 \(K^{*}/U(A)\)；对 \(x \in K\)，元素 \(v(x)\)（\(\Gamma_{\infty}\) 中的）有时称为 x 关于 v 的赋值或阶。K 上两个赋值 \(v, v'\) 若有相同的环，则称为等价的。

命题 3. 为使两个赋值 \(v, v'\) 在 \(a\)（不必交换的）域 \(K\) 上等价，必须且只需存在同构 \(\lambda\)，它把序群 \(v(K^*)\) 映到序群 \(v'(K^*)\) 上，使 \(v' = \lambda \circ v\)。

设 \(v\) 与 \(v'\) 等价；按假设，赋值 \(v\) 的环 A 与 \(v'\) 的环相同，\(v\) 与 \(v'\)（限制在 K* 上）分解为同态 \(K^* \to K^*/U(A) \xrightarrow{\mu} v(K^*)\) 与 \(K^* \to K^*/U(A) \xrightarrow{\nu} v'(K^*)\)，其中 \(\mu\) 与 \(\nu\) 是同构；此外，\(v(K^*)\)（相应地，\(v'(K^*)\)）的正元素所成之集，是在 \(\mu\)（相应地，\(\nu\)）之下 m(A) 的 \(\neq 0\) 的元素模 U(A) 的类所成之集的像；我们断定 \(\lambda = \nu \circ \overline{\mu}^{-1}\) 解决问题，而逆命题是明显的。

现在设 K 是交换域；这时，对 K 上的每个赋值 v，赋值 v 的环 A 都是 § 1 第 2 节定义 2 意义下 K 的赋值环（这证明了所用术语的合理性）；这由命题 2 (c) 与 § 1 第 2 节定理 1 (c) 立即得出。反之，回忆对每个以 K 为分式域的整环 B，整除关系 x\textbar y（等价于 \(y \in Bx\)）使 \(K^{*}\) 成为一个预序群，其相伴序群 \(\Gamma_{B}\) 是商群 \(K^{*}/U(B)\)，即 \(K^{*}\) 对 B 的可逆元素所成的群 \(U(B)\) 的商，这个群的正元素是 \(B^{*}/U(B)\) 的元素（其中 \(B^{*} = B - \{0\}\)）；映射 \(x \mapsto Bx\) 通过取商定义了从序群 \(K^{*}/U(B)\) 到 K 的非零主分式理想所成的群（以关系 \(\supset\) 排序）上的同构（《代数》第六章 § 1 第 5 节）。以 K 为分式域且使群 \(\Gamma_{A} = K^{*}/U(A)\) 全序的那些环 A 恰是 K 的赋值环（§ 1 第 2 节定理 1 (d)）。若 \(v_{A}\) 表示从 \(K^{*}\) 到 \(\Gamma_{A}\) 上的典范同态，则立即可见 \(v_{A}\)（用 \(v_{A}(0) = +\infty\) 扩张后）是 K 上的一个（称为典范的）赋值，其环为 A；每个等价于 \(v_{A}\) 的赋值都可写成 \(v = \sigma \circ v_{A}\)，其中 \(\sigma\) 是从 \(\Gamma_{A}\) 到 v 取值的群的一个子群上的同构（命题 3）；\(\sigma \circ v_{A}\) 称为 v 的典范分解。

命题 4. 设 C 是整环，K 是它的分式域，\(C^{*} = C - \{0\}\)，\(v: C \to \Gamma_{\infty}\) 是 C 上的赋值。则存在唯一的延拓 v 的 K 上的赋值 w，且 \(w(K^{*})\) 是 \(\Gamma\) 的子群，由 \(v(C^{*})\) 所生成。

由《代数》第一章 § 2 第 7 节定理 2，存在唯一的同态 \(w\)（从 \(\mathbf{K}^*\) 到 \(\Gamma\)），它延拓 \(v \mid \mathbf{C}^*\)，且 \(w(\mathbf{K}^*)\) 由 \(v(\mathbf{C}^*)\) 生成。剩下要证明 \(w\) 满足公理 \((\mathrm{VL}_{\mathrm{II}})\)。于是设 \(x' \in \mathbf{K}^*\)、\(y \in \mathbf{K}^*\) 满足 \(x + y \in \mathbf{K}^*\)；存在 \(a \in \mathbf{C}^*\) 使 \(ax \in \mathbf{C}^*\) 且 \(ay \in \mathbf{C}^*\)，从而 \(a(x + y) \in \mathbf{C}^*\)。由于 \(w\) 在 \(\mathbf{C}^*\) 上的限制满足 \((\mathrm{VL}_{\mathrm{II}})\)，

\[
w (a (x + y)) \geqslant \inf (w (a x), w (a y)).
\]

从两端消去 \(w(a)\)，我们得到

\[
w (x + y) \geqslant \inf (w (x), w (y)).
\]

